

Orbiter can create translation planes from spreads. 
The construction of translation planes from spreads 
is due to Andr\'e and Bruck, Bose (cf.~\cite{Andre54,BruckBose64}). 
In order to perform the construction, we need a field $F=\bbF_q$ 
which is the kernel of the plane, 
a spread of $k$-subspaces, 
and the groups $\PGGL(2k,F)$ and $\PGGL(2k+1,F)$.

\bigskip

Table~\ref{tab:translationplane:activity} lists Orbiter activities for translation planes.
\orbitertableofcommands{translation_plane_activity}


For instance, the command 

\sourcecode{create_translation_plane_9b}

creates the (projective) Hall plane of order 9 from the Hall spread.
In this example, we use the fact that $\PGGL(n,q) = \PGL(n,q)$ if $q$ is prime.
The example also creates a bitmap drawing of the incidence matrix of the plane, 
shown below:
\orbiteroutput{GRAPHICS/WRAPPERS/plane_catalogue_q3_k2_1_incma_draw}

\bigskip

In the next example, we will create a translation plane 
of order 16 with kernel of order 4:

\sourcecode{create_translation_plane_16_4_0}

This plane is the Hall plane, and the spread is the Hall spread. 
The stabilizer of the spread has order 1200.

\bigskip

In the next example, we will create 
a translation plane of order 16 with kernel of order 2:

\sourcecode{create_translation_plane_16_2_0}

The spread has a stabilizer of order 1008, 
which means that the associated translation plane 
has a stabilizer of order $1008 \cdot 256 = 258045.$
According to~\cite{Moorhouse:tables:16}, there are two planes 
whose associated spreads have an automorphism group of this order. 
They can be distinguished by the 2-rank of their incidence matrices. 
The Johnson-Walker plane has a 2-rank of 100. 
The Lorimer-Rahilly plane has a 2-rank of 106. 
Using Orbiter, we compute the 2-rank of the translation plane that we have created:

\sourcecode{RREF_plane_16_2_0_rank_of_incma}

It turns out that the 2-rank of our plane is 106, so the plane is Lorimer-Rahilly.


\bigskip

Let us investigate the Rao / Rao plane from Section~\ref{sec:spreads}, 
which we know is isomorphic to the spread in the Orbiter catalogue 
with number 0 and projective stabilizer of order 84.
The command  

\sourcecode{create_translation_plane_27_Rao_Rao}

creates the translation plane from the spread.
To compute the 3-rank of the incidence matrix, we issue the following comand:

\sourcecode{RREF_Rao_Rao_plane_incma_rank}

The 3-rank turns out to be 271. 
According to the Moorhouse tables~\cite{Moorhouse:tables:27}, 
the plane has Moorhouse number IV.

\bigskip

Suppose we want to study the orbits of the group of a translation plane on points and blocks. 
The following command creates the translation plane of order 27 with Orbiter number $4.$
Once done, it creates the action on points and blocks and 
invokes the Schreier orbit algorithm (see Section~\ref{sec:orbit:algorithms}).

\sourcecode{translation_plane_27_4_orbits}

The group of the plane has order $1592136.$
The degree of the action is $1514.$
The orbit lengths are
$$
1, \; 28, \; 729, \; 756.
$$

